\documentclass[11pt]{article}
\usepackage{amsmath,amssymb,amsthm,hyperref}
\newtheorem{theorem}{Theorem}
\newtheorem{lemma}[theorem]{Lemma}
\title{Permutation Averaging and the Cost of Separated Coefficients}
\author{}
\date{November 23, 2026}
\begin{document}
\maketitle

\paragraph{Scope.}
Permutation averaging gives exact positive reciprocal identities with
coefficients outside a single additive lattice. This note audits that
construction. It does not prove or disprove the universal head bound.
The permutation-sum identity and the abstract subset-moment support
argument are Lean formalized. The numerical mass evaluation, finite
separation estimate, and coefficient-pool count are manual proofs.

\section{Exact averaging over arbitrary positive values}

Let $v_0,\ldots,v_n>0$ and, for a permutation $\sigma$ of
$\{0,\ldots,n\}$, put
\[
 K_\sigma(x)=\frac{v_{\sigma(0)}}
                  {v_{\sigma(0)}+\sum_{i=1}^n v_{\sigma(i)}x_i}.
 \tag{1}
\]
The denominators have strictly positive slopes. For every Boolean input,
\[
        \frac1{(n+1)!}\sum_\sigma K_\sigma(x)
                            =\frac1{1+|x|}.
 \tag{2}
\]
Indeed let $A=\{0\}\cup\{i:x_i=1\}$. For each $i\in A$, the
sum over permutations of $v_{\sigma(i)}/\sum_{j\in A}v_{\sigma(j)}$
is the same, by composing permutations with a transposition inside $A$.
Summing these equal quantities over $i\in A$ gives $(n+1)!$,
which proves (2). This argument is kernel checked by
\texttt{HeadComplexity.permutation\_reciprocal\_ratio\_sums\_eq}
and \texttt{HeadComplexity.permutation\_reciprocal\_ratio\_sum}
in \texttt{Results/PermutationReciprocalAverage.lean}.

If several $v_i$ are equal, permutations with identical coefficient
vectors can be merged. The finite-capacity identities in
\texttt{squarefree-reciprocal-basis.tex} result by using two distinct
values and, more generally, grouping tokens into input colors.
Merging arbitrary widely separated values behaves differently.

\section{The anchored subset-distribution invariant}

Let a random permutation have an arbitrary probability distribution,
not necessarily uniform, and write
\[
 R=\{i\in[n]:\sigma(i)<\sigma(0)\},\qquad
 m_T=\Pr(T\subseteq R),\quad T\subseteq[n].
\]
For each $S\subseteq[n]$, inclusion-exclusion gives
\[
 \Pr(R=S)=\sum_{U\subseteq[n]\setminus S}(-1)^{|U|}m_{S\cup U}.
 \tag{3}
\]
If $m_T=1/(|T|+1)$ for every $T$, then
\[
 \Pr(R=S)=\int_0^1t^{|S|}(1-t)^{n-|S|}\,dt
                 =\frac1{(n+1)\binom n{|S|}}>0.
 \tag{4}
\]
The first equality follows by expanding $(1-t)^{n-|S|}$ and
integrating each monomial. The second follows by repeated integration
by parts, or directly by counting the permutations whose indices
before the distinguished index are exactly $S$.
Thus even this one distinguished-index version requires at least
$2^n$ permutations in its support. It is sharp as a distributional
statement: choose one permutation for each $S$ and give it the weight
in (4). The moments then equal $1/(|T|+1)$ by the same integral.

The upper-subset moment transform is injective, as checked by
\texttt{HeadComplexity.subset\_moment\_zero} and
\texttt{HeadComplexity.subset\_moment\_ext}. The declaration
\texttt{HeadComplexity.card\_le\_of\_positive\_subset\_moments}
in \texttt{Results/SubsetMomentRigidity.lean} proves that any finite
signed family with these moments needs at least $2^n$ atoms once
strict positivity of the target masses is supplied. Formula (4) is
proved manually above; it is not assumed to have been formalized.

The inclusion-exclusion method is related to the proof of Theorem 7
in Broder, Charikar, Frieze, and Mitzenmacher,
\emph{Min-Wise Independent Permutations}, JCSS 60 (2000),
\href{https://www.cse.iitb.ac.in/~soumen/readings/papers/BroderCFM2000minwise.pdf}
{Section 3.4}. Their stronger requirement treats every distinguished
index and gives a larger support lower bound. The following application
to positive reciprocal kernels needs only (3)--(4), which were proved
above.

\begin{theorem}[Separated permutation averages need a full cube of terms]
Let $n\ge1$ and $q>1$. Set $v_j=q^j$ for $0\le j\le n$ and
\[
             M_n=(n+1)\binom n{\lfloor n/2\rfloor}.
\]
Suppose permutations $\sigma_1,\ldots,\sigma_H$ and nonnegative
weights $p_1,\ldots,p_H$ satisfy the exact identity
\[
           \sum_{h=1}^H p_h K_{\sigma_h}(x)=\frac1{1+|x|}
                   \quad(x\in\{0,1\}^n).
 \tag{5}
\]
If $q>n2^n M_n$, then $H\ge2^n$. Repeated permutations count as
separate occurrences and cannot improve this lower bound.
\end{theorem}
\begin{proof}
The input zero in (5) gives $\sum_h p_h=1$. For an input with
support $T$, the kernel in (1) differs by at most $n/q$ from
$\mathbf1_{\{T\subseteq R_\sigma\}}$. If the indicator is one,
all other active powers are at most $q^{\sigma(0)-1}$, giving a
kernel value at least $1/(1+n/q)\ge1-n/q$. If it is zero,
one active power is at least $q^{\sigma(0)+1}$, giving a kernel
value at most $1/q$.

Averaging and using (5) proves
$|m_T-1/(|T|+1)|\le n/q$ for every $T$. In (3) this gives
\[
 \left|\Pr(R=S)-\frac1{(n+1)\binom n{|S|}}\right|
                 \le\frac{n2^{n-|S|}}q\le\frac{n2^n}q.
 \tag{6}
\]
The prescribed inequality makes the right side smaller than $1/M_n$,
which is no larger than the positive central value in (6). Hence every
$S$ has positive probability. Each permutation has only one set $R$,
so at least $2^n$ permutations are required.
\end{proof}

The support bound is sharp for positive exact averaging at these
separated values. The full permutation average (2) is feasible, and
all its feature vectors have coordinate one at $x=0$. If a positive
representation uses more than $2^n$ vectors, they are linearly
dependent. A dependence has coefficient sum zero by the anchor
coordinate. Move the weights along it until at least one reaches
zero, preserving the table and nonnegativity. Repeating gives a
representation with at most $2^n$ terms. Thus the minimum is exactly
$2^n$ under the theorem's separation hypothesis. For rational $q$,
the elimination can be performed over the rationals.

The same computation gives a cancellation requirement for signed
weights. Put $V=\sum_h|p_h|$. The error in (6) is then at most
$n2^n V/q$. If $H<2^n$, some subset has no supporting permutation,
so its total signed mass is zero. Therefore every signed identity
(5) with $H<2^n$ must satisfy
\[
                         V\ge\frac{q}{n2^n M_n}.
 \tag{7}
\]
This is a bound for the specified permutation construction. The
manuscript permits unrestricted signed weights, so (7) is not a
lower bound for general head complexity.

\section{A common small coefficient pool cannot be universal}

There is a separate obstruction that does not require an averaging
identity, separated values, or positive numerator weights. Suppose
each denominator, after a positive scalar normalization absorbed into
its numerator, has all $n+1$ affine coefficients in a common pool
$v_1,\ldots,v_M$. The pool values, assignments, numerators, and bias
may all depend on the target. Only positivity of the denominators
on the cube is imposed; allowing mixed slopes enlarges the family.

\begin{theorem}[Coefficient-pool count]
Put $N=2^n$, $p_1=H(n+1)+1$, and $p=p_1+M\le N$.
The number of strict sign tables in this family is at most
\[
 M^{H(n+1)}
 \left(\frac{32768N}{p}\right)^p
                     2^{p_1}\max\{2,H\}^{M}.
 \tag{8}
\]
In particular, for fixed $C>0$, the family is not universal if
$H\le C2^n/n^2$ and $\log_2 M=o(n)$.
\end{theorem}
\begin{proof}
There are at most $M^{H(n+1)}$ coefficient assignments. For one
assignment clear the positive product of the $H$ denominators.
At each vertex the resulting polynomial is of degree one in the
$p_1$ numerator and bias parameters and of degree at most $H$
in the $M$ pool values. The two-block sign bound recorded in
\texttt{mechanism-audit.tex} proves (8). Any restrictions on signs,
coincidences, or assignments can only reduce the count.

At the stated budget, $p_1=O_C(N/n)$ and $M=2^{o(n)}$.
The logarithm of the assignment factor is at most
$O_C((N/n)\log M)=o(N)$. The sign-bound contribution from $p_1$
is $O_C(N\log n/n)$, and the additional contributions from the
pool parameters are at most $O(Mn)=o(N)$. Thus the total is $o(N)$.
There are $2^N$ Boolean sign tables, proving nonuniversality.
\end{proof}

All permutations of a single target-dependent vector of $n+1$
coefficients belong to this pool family. Thus varying their signed
numerator weights does not remove this particular obstruction. An
unrestricted collection of independently adaptive denominators is
not included in it.

\end{document}
